% %%
\documentclass[lettersize,journal]{IEEEtran}
%\documentclass[lettersize,journal,draft]{IEEEtran}

\usepackage{amsmath,amsfonts}
\usepackage{algorithmic}
\usepackage{algorithm}
\usepackage{array}
\usepackage[caption=false,font=normalsize,labelfont=sf,textfont=sf]{subfig}
\usepackage{textcomp}
\usepackage{stfloats}
\usepackage{url}
\usepackage{verbatim}
\usepackage{graphicx}
\usepackage{cite}
\hyphenation{op-tical net-works semi-conduc-tor IEEE-Xplore}
\def\BibTeX{{\rm B\kern-.05em{\sc i\kern-.025em b}\kern-.08em
    T\kern-.1667em\lower.7ex\hbox{E}\kern-.125emX}}
\usepackage{balance}
\usepackage{xcolor}
\newtheorem{proposition}{Proposition}
% Revision-round number markers (2026-11-02). Every marked number is listed in
% Queued_Edits/revision_numbers_ledger.md.
\newcommand{\syn}[1]{\textcolor{red}{#1}}            % placeholder data, replaced by the campaign
\newcommand{\rev}[1]{#1}                             % simulation or analysis result, accepted 2026-11-03 (black); rerun if campaign settings change
\newcommand{\open}[1]{\textcolor{magenta}{[#1]}}     % protocol setting not yet decided
\newcommand{\newnum}[1]{\textcolor{blue}{#1}}        % simulation result added 2026-11-04, not yet reviewed (rename to \rev once approved)
% For submission: redefine all three as {#1} (and \open as {}) after the placeholders are gone.
\begin{document}
\title{Adaptive Navigation of Multirobot Systems to Critical Points in 2D Vector Fields in Simulation and Experiment}
\author{Christopher Waight and Christopher A. Kitts,~\IEEEmembership{Senior Member,~IEEE}
\thanks{Manuscript received Month Day, Year; revised Month Day, Year.}
\thanks{C. Waight is a PhD Candidate in the Robotic Systems Laboratory, Department of Mechanical Engineering, Santa Clara University, Santa Clara, CA 95053 USA (e-mail: cwaight@scu.edu)}
\thanks{C. A. Kitts is Professor and Director of the Robotic Systems Laboratory, Department of Mechanical Engineering, Santa Clara University, Santa Clara, CA 95053 USA (e-mail: ckitts@scu.edu)}
}
\markboth{Journal of \LaTeX\ Class Files,~Vol.~XX, No.~X, Month~Year}%
{Waight and Kitts: Adaptive Navigation of Multirobot Systems to Critical Points in 2D Vector Fields}

%\IEEEpubid{0000--0000/00\$00.00~\copyright~2025 IEEE}

\maketitle

\begin{abstract}
    There are a variety of interesting phenomena that exist at the critical points of environmental vector fields. Existing approaches for navigation to critical points either require prior knowledge of the environment or use local information to provide directional guidance only. This paper presents a multirobot method for estimating critical point locations without prior field knowledge, using only real-time local measurements from a minimum of three robots. The sensing is spatially distributed across the robots, while the estimate and control laws are computed on a central computer.
    This estimation enables new control laws to navigate to these points and to orbit them at a fixed radius.
    The estimation error is bounded, and the attraction error stays bounded under actuator lag, sampling, and friction.
    Simulation verifies both attraction and orbital control laws across six canonical linear vector fields and three nonlinear fields using three omnidirectional robots. Hardware validation within a lab-based experimental testbed demonstrates a \syn{100}\% (saddle) and \syn{86}\% (vortex) attraction convergence success rate within 0.05~m across \syn{158} trials, with mean position errors of \syn{0.037}~m or less from the true critical points. An additional \open{36} orbital trials at \open{six} commanded radii held bounded orbits, with radial bias within a quarter of the radius from \syn{0.40}~m (vortex) and \syn{0.50}~m (saddle) upward. % [SYNTHETIC] hardware numbers; [OPEN] orbit trial count and radii
    These findings establish this estimation method under lab conditions, where printed fields do not advect the robots. The present controller does not yet compensate for oceanic and atmospheric flows that push the platform.

\end{abstract}


\begin{IEEEkeywords}
Multirobot systems, adaptive navigation, vector fields, critical points, cooperative estimation, cluster control
\end{IEEEkeywords}


\section{Introduction}

\IEEEPARstart{M}{ultirobot} systems enhance redundancy, throughput, coverage, and reconfigurability of robotic systems in applications ranging from warehouse fulfillment to search and rescue \cite{1}. Equipped with instruments, they form mobile sensor networks that exploit their spatial distribution to estimate and characterize environmental features across a region \cite{2,3,4}. Adaptive navigation (AN) techniques use these estimates to direct the network to or along features of interest. Unlike survey-oriented navigation, which exhaustively traverses a field and identifies features afterward, AN quickly locates local features and follows them in dynamic environments \cite{5,6,7}.

AN in scalar fields is well explored. Each robot collects a single scalar measurement as the input to a navigation policy that reaches and follows features such as extrema, saddle points, ridges, trenches, contour lines, and fronts \cite{5,6,7,8,9}. Strategies include gradient following \cite{10}, Hessian estimation \cite{4}, extended information consensus filters \cite{11,13}, constrained nonsmooth distributed optimization \cite{14}, and bio-inspired search \cite{15}, and several such systems have been field-deployed on ground, surface, and aerial vehicles for environmental monitoring, pollution tracking, and search and rescue \cite{10,18,19,20}.

By contrast, only a few works address navigation of physically sensed vector fields \cite{21,22,23}. Much of the broader vector field literature constructs artificial fields for path-following and obstacle avoidance, a separate problem from sensing a real flow \cite{24,25,26,28}, and field campaigns that track features carried by an actual ocean current underscore the distinction \cite{29}.

The oceanic and atmospheric environments that motivate this work are governed by a spatially varying velocity or force, and their transport and accumulation behavior is organized by the critical points of the field, the coordinates where the flow vanishes \cite{32,36}. A surface convergence zone acts as a sink, gathering buoyant debris and drifting bodies; this makes such zones high probability search regions after a loss at sea \cite{37,38}. Containing a ruptured wellhead or a harbor discharge requires reaching the source against its own outward flow \cite{40,41}. A mesoscale ocean eddy serves as a center, trapping a body of water, along with the plankton bloom inside it, and carrying both across a basin for weeks \cite{43,44,45}. Monitoring the bloom over its lifetime therefore requires a sensor to orbit the drifting core at a fixed standoff; a ground-fixed station would not stay with the eddy. These applications require the coordinate of the feature itself, whether to drive the multirobot team onto it or to hold it at a fixed standoff radius; a heading toward the feature is not sufficient.

Work that senses the physical field falls into two camps. One camp presumes prior knowledge of the environment, which is unavailable in the unmapped settings that motivate this work \cite{30,31}. The other operates online in two behaviors: minimally actuated drifters that ride and station-keep along ambient streamlines, and PIM-triple-inspired teams that straddle a saddle and track a single stable or unstable manifold from local velocity measurements \cite{23,32,33}. Both supply a heading along the flow or track the manifold of a single assumed saddle; neither recovers the coordinate of a critical point together with its identity across the full taxonomy of feature types. Estimating the locations of critical points in vector fields remains an open problem for multirobot systems.

Recovering those coordinates is harder than following a heading, for reasons of both information and implementation. The information problem is that the location and classification of a critical point are absent from any single velocity sample; they are carried by the spatial derivatives of the field \cite{46}. A single platform must move, sample again, and difference its readings, which conflates the spatial gradient with the field's change in time and can return features that are not physically present. A distributed team removes this temporal aliasing by sampling at one instant \cite{22,23,47}, but recovering the derivatives is still not recovering the feature. Magnitude-based methods \cite{48} reduce the flow to a scalar and yield a descent direction, but the magnitude vanishes at every critical point regardless of type, and estimating the two velocity components as independent scalar fields still does not locate the point. Both the coordinate and the identity are properties of the assembled Jacobian. Inverting it recovers the location where the field vanishes, and its eigenvalues determine the type \cite{36}.

The implementation problem is that the local field must be reconstructed on physical hardware from noisy distributed measurements. The conditioning of this reconstruction degrades as the formation approaches collinearity, so the sampling shape must be held by a formation controller under real actuator dynamics and communication latency. Both the estimator and the mechatronics of the cluster limit the accuracy of the recovered features.

This paper introduces a cooperative estimation strategy that recovers the local linear structure of the field. Using a minimum of three robots, the cluster estimates the local Jacobian in a single time step and treats critical point navigation as a local-to-global inference. Inverting the Jacobian predicts the location of the nearby critical point that organizes the surrounding flow, and its eigenvalues indicate the type, with no global map required. The same estimate supports two control laws, a proportional law that drives the cluster centroid onto the critical point, and an orbital law that holds the centroid at a commanded standoff radius. In the present implementation the robots sense and move, and a central computer forms the estimate and computes the commands.

This paper makes four contributions. First, it presents an algorithm that estimates critical point locations from instantaneous, spatially distributed measurements alone. It also shows how field curvature, formation size, and measurement error set the error of the estimate. Second, it proves that three non-collinear robots are the minimum number required for this estimation, and that they suffice for every nondegenerate affine field. Third, it demonstrates that this estimate suffices to drive navigation to and around critical points across six vector field types and three nonlinear fields. It also states conditions under which the attraction error stays bounded despite actuator lag, sampled updates, and friction. Fourth, it validates the cooperative estimation and resulting navigation in \open{campaign} hardware experiments that demonstrate repeatable critical point estimation, attraction, and orbit control. % [OPEN] campaign trial count; the submitted 169 is superseded



\section{Critical Point Detection and Control}

\subsection{Vector Field Environments and Critical Points}

In its most general form, a vector field assigns a vector $\mathbf{v}(\mathbf{p}) \in \mathbb{R}^m$ to each point $\mathbf{p} \in \mathbb{R}^n$ within a given domain, where $m \geq 2$, making vector fields well suited to represent directional physical quantities such as fluid velocity, electromagnetic force, or gravitational force. For the planar fields analyzed in this work, this takes the specific form of a 2D vector $\mathbf{v}(\mathbf{p}) = (u(\mathbf{p}), v(\mathbf{p}))$ assigned to each coordinate $\mathbf{p} = (x, y)$.

Unlike scalar fields where $m=1$, vector fields can exhibit features such as critical points $\mathbf{p}^*$ where the field vanishes, $\mathbf{v}(\mathbf{p}^*) = \mathbf{0}$. Critical points model phenomena such as vortices arising from tidal flows over submarine caves, or plumes when searching for pollution sources \cite{49}.

Near a critical point, the field is well approximated by its linearization, and the local topology is determined by the Jacobian matrix $\mathbf{J} = \nabla\mathbf{v}|_{\mathbf{p}}$. The eigenvalues of $\mathbf{J}$ evaluated at $\mathbf{p}^*$ determine the topological type. For $2 \times 2$ systems with nonzero determinant, six standard critical point types arise from the eigenvalue structure, as shown in Fig.~\ref{fig:allfield} and classified in Table~\ref{tab:field_summary}.

\begin{figure}[!t]
    \centering
    \includegraphics[width=0.95\columnwidth]{six_vector_fields_3x2.png}
    \caption{Six linear vector field types with critical points.}
    \label{fig:allfield}
\end{figure}

Eigenvalue analysis can also identify degenerate scenarios such as cases where the determinant is zero. In uniform flow fields, for instance, no critical point exists. While these scenarios are detectable through the estimation framework of Section~II-B, they fall outside the scope of the present work. A nonlinear field can also hold several critical points, separated by curves on which the determinant is zero.

\begin{table}[!t]
    \centering
    \caption{Classification of critical points by Jacobian eigenvalue structure ($\lambda$ = eigenvalue; a complex pair $\alpha \pm i\omega$ has real part $\alpha$, imaginary part $\omega$).}
    \label{tab:field_summary}
    \footnotesize
    \setlength{\tabcolsep}{2pt}
    \begin{tabular}{lll}
    \hline
    \textbf{Field Type} & \textbf{Critical Point} & \textbf{Eigenvalues} \\
    \hline
    Vortex & Center & $\pm i\omega$ \\
    Sinking Vortex & Stable Spiral & $\alpha \pm i\omega$, $\alpha < 0$ \\
    Sink & Stable Node & $\lambda_1, \lambda_2 < 0$ \\
    Source & Unstable Node & $\lambda_1, \lambda_2 > 0$ \\
    Spewing Vortex & Unstable Spiral & $\alpha \pm i\omega$, $\alpha > 0$ \\
    Saddle & Saddle Point & $\lambda_1 < 0 < \lambda_2$ \\
    \hline
    \end{tabular}
\end{table}


\subsection{Cooperative Critical Point Estimation}
\noindent

We refer to the group of robots as a cluster: a formation treated as a single virtual rigid body whose centroid, orientation, and internal geometry can be independently specified and controlled \cite{50,51}.

Given a cluster of three robots sampling a planar vector field with robot positions $\mathbf{p}_i = (x_i, y_i)$ and corresponding vector field measurements $(u_i, v_i)$ for $i \in \{1, 2, 3\}$, we can approximate the vector field components across the cluster's spatial extent using linear planar models:
\begin{equation}\label{eq:2}
\hat{U}(x,y) = ax + by + c = u(x,y) + \mathcal{O}(\ell^2)
\end{equation}
\begin{equation}\label{eq:3}
\hat{V}(x,y) = dx + ey + f = v(x,y) + \mathcal{O}(\ell^2)
\end{equation}
where hats denote quantities estimated from the measurements and $\ell$ is the size of the region spanned by the robots and the point $(x,y)$. The $\mathcal{O}(\ell^2)$ term carries the field's curvature, the second-order Taylor terms that a plane cannot represent.

The coefficients are determined by solving the linear systems:
\begin{equation}\label{eq:4}
\mathbf{A}\hat{\boldsymbol{\theta}}_u = \mathbf{z}_u, \quad \mathbf{A}\hat{\boldsymbol{\theta}}_v = \mathbf{z}_v
\end{equation}
where:
\begin{equation}\label{eq:5}
\mathbf{A} = \begin{bmatrix} x_1 & y_1 & 1 \\ x_2 & y_2 & 1 \\ x_3 & y_3 & 1 \end{bmatrix}, \quad \hat{\boldsymbol{\theta}}_u = \begin{bmatrix} a \\ b \\ c \end{bmatrix}, \quad \hat{\boldsymbol{\theta}}_v = \begin{bmatrix} d \\ e \\ f \end{bmatrix}
\end{equation}
and $\mathbf{z}_u = [u_1, u_2, u_3]^T$, $\mathbf{z}_v = [v_1, v_2, v_3]^T$ are the collected field measurements.

The estimated critical point $\hat{\mathbf{p}}^* = (\hat{x}^*, \hat{y}^*)$ satisfies $\hat{U}(\hat{x}^*, \hat{y}^*) = 0$ and $\hat{V}(\hat{x}^*, \hat{y}^*) = 0$, which in matrix form reads $\hat{\mathbf{J}}\hat{\mathbf{p}}^* = -\hat{\mathbf{h}}$, where the estimated Jacobian matrix and offset vector are:
\begin{equation}\label{eq:8}
\hat{\mathbf{J}} = \begin{bmatrix} a & b \\ d & e \end{bmatrix}, \quad \hat{\mathbf{h}} = \begin{bmatrix} c \\ f \end{bmatrix}
\end{equation}

Thus, provided $\det(\hat{\mathbf{J}}) \neq 0$:
\begin{equation}\label{eq:9}
\hat{\mathbf{p}}^* = -\hat{\mathbf{J}}^{-1}\hat{\mathbf{h}}
\end{equation}

The eigenvalues of $\hat{\mathbf{J}}$ can be used to identify the type of nearby critical point (Table~\ref{tab:field_summary}).

Equation (\ref{eq:9}) returns the zero of the fitted affine field $\hat{\mathbf{v}}(\mathbf{p}) = \hat{\mathbf{J}}\mathbf{p} + \hat{\mathbf{h}} = \hat{\mathbf{J}}(\mathbf{p} - \hat{\mathbf{p}}^*)$, which need not be a zero of the true field. Evaluating it at any critical point $\mathbf{p}^*$ of the true field gives
\begin{equation}\label{eq:scope}
\hat{\mathbf{p}}^* - \mathbf{p}^* = -\hat{\mathbf{J}}^{-1}\hat{\mathbf{v}}(\mathbf{p}^*).
\end{equation}
The true field vanishes at $\mathbf{p}^*$, so $\hat{\mathbf{v}}(\mathbf{p}^*)$ is the error of the affine model there. If the true field is affine over a region containing both the cluster and $\mathbf{p}^*$, the models (\ref{eq:2})--(\ref{eq:3}) reproduce it without error and the recovery is exact at any formation size. On a smooth nonlinear field, the fit carries a first-order linearization error that scales with the local curvature of the field; this error decreases as the formation contracts toward the critical point. The inference from $\hat{\mathbf{p}}^*$ to a true critical point is therefore valid only where the affine model remains accurate from the cluster out to $\mathbf{p}^*$.

The errors in $\hat{\mathbf{J}}$ and $\hat{\mathbf{h}}$ come from the same measurements and are correlated, but (\ref{eq:scope}) carries them jointly through $\hat{\mathbf{v}}(\mathbf{p}^*)$. Taking norms gives $\|\hat{\mathbf{p}}^* - \mathbf{p}^*\| \le \|\hat{\mathbf{v}}(\mathbf{p}^*)\|/\sigma_{\min}(\hat{\mathbf{J}})$. The singular value $\sigma_{\min}(\hat{\mathbf{J}})$ is the slowest rate at which the fitted field changes in any direction, so a small value amplifies the model error.

Let $\mathbf{p}_c = \frac{1}{3}\sum_i\mathbf{p}_i$ be the cluster centroid, $\mathbf{e} = \mathbf{p}_c - \mathbf{p}^*$ its displacement from the critical point, and $R$ the distance from $\mathbf{p}_c$ to each robot of the equilateral formation. Let $M$ bound the Hessian norms of $u$ and $v$ over the region holding the robots and $\mathbf{p}^*$, and let $\bar{\varepsilon}$ bound each robot's measurement error. Each reading then lies within $MR^2/\sqrt{2} + \bar{\varepsilon}$ of the linearization about $\mathbf{p}_c$. The fit carries these deviations to $\mathbf{p}^*$ through the barycentric weights $w_i = \frac{1}{3} - \frac{2}{3}(\mathbf{p}_i - \mathbf{p}_c)^T\mathbf{e}/R^2$, whose magnitudes sum to at most $1 + \sqrt{2}\|\mathbf{e}\|/R$. Adding the linearization's own error $M\|\mathbf{e}\|^2/\sqrt{2}$ at $\mathbf{p}^*$ and merging the curvature terms gives
\begin{equation}\label{eq:bound}
\|\hat{\mathbf{p}}^* - \mathbf{p}^*\| \le \frac{\overbrace{\tfrac{M}{\sqrt{2}}(R + \|\mathbf{e}\|)^2}^{\text{curvature}} + \overbrace{\bigl(1 + \tfrac{\sqrt{2}\|\mathbf{e}\|}{R}\bigr)\bar{\varepsilon}}^{\text{measurement error}}}{\sigma_{\min}(\hat{\mathbf{J}})}.
\end{equation}
Formation size moves the two terms in opposite directions. A large formation raises the curvature term, which grows with the square of the largest distance from a robot to $\mathbf{p}^*$. A small formation raises the measurement term, which grows with the extrapolation distance $\|\mathbf{e}\|/R$ measured in formation sizes. The bound assumes a true critical point exists nearby and does not establish that it does.

Each robot carries its own sensor, so measurement noise is independent across robots. For zero-mean noise of covariance $\sigma_v^2\mathbf{I}_2$ at each robot on a locally affine field, linearizing (\ref{eq:scope}) in the noise gives
\begin{equation}\label{eq:cov}
\operatorname{Cov}(\hat{\mathbf{p}}^*) \approx \sigma_v^2\, c_w\, \mathbf{J}^{-1}\mathbf{J}^{-T},
\end{equation}
where $c_w = \sum_i w_i^2 = \frac{1}{3} + \frac{2}{3}\|\mathbf{e}\|^2/R^2$ for an equilateral formation. As the cluster reaches the critical point, $c_w \to 1/3$, the variance factor of the mean of three independent readings.


\subsection{Minimum Robot Requirements}

Equations (\ref{eq:2})--(\ref{eq:3}) each contain three coefficients. From (\ref{eq:4}), each robot contributes one equation to each system: robot $i$ provides $u_i = ax_i + by_i + c$ for the U-component and $v_i = dx_i + ey_i + f$ for the V-component. To uniquely determine three unknowns requires at least three equations. Three robots are therefore the minimum required for critical point estimation.


Three non-collinear robots are also sufficient: $\det(\mathbf{A})$ equals twice the signed area of the formation triangle, so $\mathbf{A}$ is invertible exactly when the robots are not collinear.

\begin{proposition}[Minimality]\label{prop:min}
For an affine field $\mathbf{v}(\mathbf{p}) = \mathbf{J}(\mathbf{p}-\mathbf{p}^*)$ with $\det(\mathbf{J}) \neq 0$, locating $\mathbf{p}^*$ of every such field from instantaneous measurements requires at least three non-collinear robots, and three suffice.
\end{proposition}

\begin{IEEEproof}
Any $n \le 2$ robots, or any collinear robots, lie on a line $L = \{\mathbf{p}_0 + s\mathbf{d}\}$ and measure $\mathbf{m}_0 + s_i\mathbf{m}_1$, with $\mathbf{m}_0 = \mathbf{J}(\mathbf{p}_0 - \mathbf{p}^*)$ and $\mathbf{m}_1 = \mathbf{J}\mathbf{d}$. Take a field with $\mathbf{p}^* \notin L$. For any $\mathbf{q} \notin L$, the field $\mathbf{J}_q(\mathbf{p} - \mathbf{q})$ with $\mathbf{J}_q = [\,\mathbf{m}_0\ \ \mathbf{m}_1\,][\,\mathbf{p}_0 - \mathbf{q}\ \ \mathbf{d}\,]^{-1}$ is nonsingular and reproduces every measurement on $L$, so the measurements cannot determine $\mathbf{p}^*$. Its determinant $\det(\mathbf{J}_q) = \det(\mathbf{J})\det[\,\mathbf{p}_0 - \mathbf{p}^*\ \ \mathbf{d}\,]/\det[\,\mathbf{p}_0 - \mathbf{q}\ \ \mathbf{d}\,]$ changes sign as $\mathbf{q}$ crosses $L$, so the type is also undetermined. Conversely, for three non-collinear robots $\mathbf{A}$ is invertible, so $\hat{\mathbf{J}} = \mathbf{J}$ and $\hat{\mathbf{p}}^* = \mathbf{p}^*$, and the same holds for least squares with $n > 3$ robots that are not all collinear.
\end{IEEEproof}


\subsection{Adaptive Navigation Control Laws}
\label{sec:control_primitives}

We develop two control primitives that drive the cluster centroid $\mathbf{p}_c$ toward or around the estimated critical point $\hat{\mathbf{p}}^*$. The mapping from cluster-level velocity commands to individual robot velocities is handled by the cluster space controller described in Section~III.

\subsubsection{Critical Point Attraction}
To navigate the cluster center $\mathbf{p}_c$ to $\hat{\mathbf{p}}^*$, we apply a proportional control law:
\begin{equation}\label{eq:10}
\dot{\mathbf{p}}_c = k(\hat{\mathbf{p}}^* - \mathbf{p}_c)
\end{equation}
where $k > 0$ is the proportional gain. With $\hat{\mathbf{e}} = \mathbf{p}_c - \hat{\mathbf{p}}^*$ and a stationary estimate, the closed-loop error dynamics are $\dot{\hat{\mathbf{e}}} = -k\hat{\mathbf{e}}$, so $\hat{\mathbf{e}}(t) = \hat{\mathbf{e}}(0)e^{-kt}$. This guarantees global exponential convergence to the estimated critical point with time constant $1/k$, provided the estimate is stationary.

\subsubsection{Orbital Navigation}
To maintain a circular orbit of radius $r_d$ around $\hat{\mathbf{p}}^*$, we define the radial vector $\mathbf{r} = \hat{\mathbf{p}}^* - \mathbf{p}_c$ with magnitude $r = \|\mathbf{r}\|$. The unit radial vector is $\mathbf{n} = \mathbf{r}/r$, and the unit tangent vector $\mathbf{t}$ is obtained by rotating $\mathbf{n}$ by $-90^\circ$. The orbital velocity command is:

\begin{equation}\label{eq:11}
\dot{\mathbf{p}}_c = k_t \mathbf{t} + k_r(r - r_d)\mathbf{n}
\end{equation}

Here, $k_t$ determines the orbital speed (where $\text{sgn}(k_t)$ dictates direction) and $k_r > 0$ drives radial convergence.
This control law decouples radial and angular dynamics. The radial error $e_r = r - r_d$ follows $\dot{e}_r = -k_r e_r$, ensuring exponential convergence to $r_d$. As $r \to r_d$, the angular velocity stabilizes to $\omega = k_t/r_d$, producing steady circular motion.

\section{Multilayer Control Architecture}


The multirobot system control architecture chosen is an established multilayer approach successfully demonstrated in prior cluster space control studies in both 2D \cite{5,8,10} and 3D \cite{7,9} environments. Fig.~\ref{fig:control_architecture} illustrates this three-layered hierarchical structure. The architecture itself is not a contribution of this work. This paper provides the estimation function that fills the feature estimation block and the attraction and orbital controllers that fill the navigation controller block. The vector field measurements enter the architecture only at the adaptive navigation layer, which uses them to form the critical point estimate; the lower layers operate on velocity commands alone. Each layer is discussed briefly.



\subsection{Robot Controller Layer}
\noindent 

Each robot's controller receives velocity commands and translates them into motion. Each robot also samples the local vector field $\mathbf{v}(\mathbf{p}_i) = (u_i, v_i)$ and reports its global position $\mathbf{p}_i$, providing the inputs to the estimation framework described in Section~II.

Details of the hardware platform are described in Section~V. In simulation, each robot is modeled as a single omnidirectional point mass with a first-order response, Coulomb friction, and a speed limit. While the robot moves, its velocity $\mathbf{v}_{r,i} = \dot{\mathbf{p}}_i$ follows its command $\mathbf{v}_{\text{des},i}$ as
\begin{equation}
    \dot{\mathbf{v}}_{r,i}(t) = -\frac{1}{\tau}\mathbf{v}_{r,i}(t) + \frac{1}{\tau}\Bigl(\mathbf{v}_{\text{des},i}(t) - v_c\frac{\mathbf{v}_{r,i}(t)}{\|\mathbf{v}_{r,i}(t)\|}\Bigr),
    \label{eq:continuous}
\end{equation}
where $\tau$ is the actuator time constant \cite{4,7} and the last term is Coulomb friction, an extra command of magnitude $v_c$ opposing the motion. The time constant on the Decabot hardware was found to be approximately 0.28 s.  The robots are also modeled with a maximum speed of $v_{\max} = 0.4$~m/s and a minimum speed of $v_c = 0.025$~m/s needed to overcome stiction (remeasured for this work; an earlier characterization of the testbed reported 0.05~m/s \cite{22}).

For digital implementation at control period $\Delta t$ and cycle $j$, the command is held over each period and friction acts against the frictionless update $\tilde{\mathbf{v}}$, so discretization of (\ref{eq:continuous}) yields
\begin{equation}\label{eq:momentum}
\begin{split}
&\tilde{\mathbf{v}} = \alpha\,\mathbf{v}_{r,i}[j] + (1-\alpha)\,\mathbf{v}_{\text{des},i}[j],\\
&s = \min\bigl\{v_{\max},\, \max\{0,\, \|\tilde{\mathbf{v}}\| - (1-\alpha)v_c\}\bigr\},\\
&\mathbf{v}_{r,i}[j+1] = s\,\tilde{\mathbf{v}}/\|\tilde{\mathbf{v}}\|,
\end{split}
\end{equation}
where $s$ is the new speed and the momentum coefficient $\alpha = e^{-\Delta t/\tau}$. At a 10~Hz control rate, $\alpha \approx 0.70$. Friction removes $(1-\alpha)v_c$ of speed each cycle and stops the robot without reversing it, so a robot at rest moves only once its command exceeds $v_c$. The cap $v_{\max}$ applies last, and the position advances as $\mathbf{p}_i[j+1] = \mathbf{p}_i[j] + \Delta t\,\mathbf{v}_{r,i}[j+1]$. The sensed field enters only the estimate, not the robot dynamics.


\begin{figure}[!t]
    \centering
    \includegraphics[width=0.95\columnwidth]{control_architecture_2.png}
    \caption{Hierarchical control architecture: robot, cluster space, and adaptive navigation layers, as in \cite{5,8}. The contents of the feature estimation and navigation controller blocks are new in this paper.}
    \label{fig:control_architecture}
\end{figure}

\subsection{Cluster Space Controller Layer}

The cluster space controller maintains the cluster's formation geometry while executing centroid velocity commands from the adaptive navigation layer. It translates cluster-level velocity commands into individual robot velocity commands using an inverse kinematic Jacobian, and transforms robot position data back into cluster variables for feedback control. In this study, we consider a cluster of three robots forming a triangular configuration. The shape can be parameterized using Side-Angle-Side (SAS) representation with parameters $(p, \beta, q)$, where $p$ is the distance between robots 1 and 2, $\beta$ is the interior angle at robot 2, and $q$ is the distance between robots 2 and 3. Differentiating the forward and inverse kinematic equations gives analytic forms for the cluster space kinematic Jacobian and its inverse \cite{50}.


\subsection{Adaptive Navigation Layer}

The adaptive navigation layer is the highest level of the control architecture. It specifies cluster mobility and shape based on real time sensor data from each robot using three functional blocks. A feature estimation block computes the critical point estimate $\hat{\mathbf{p}}^*$ using robot positions and vector field measurements as described in Section~II-B, a navigation controller generates desired centroid velocity commands $\dot{\mathbf{p}}_c$ using the control laws described in Section~II-D, and a formation controller regulates the cluster shape variables $(p, \beta, q)$ to desired values using proportional error feedback. The navigation and formation velocity commands are combined into a complete cluster-space velocity vector, which is passed to the cluster space controller layer for transformation into individual robot commands.

\subsection{Closed-Loop Stability and Tracking Bounds}
\label{sec:stability}

Section~II-D assumes a stationary estimate and an ideal centroid, but the estimate changes every cycle and the centroid responds through the sampled lag and Coulomb friction.

\subsubsection{Sampled Loop}
At cycle $j$, write the estimate as $\hat{\mathbf{p}}^*[j] = \mathbf{p}^* + \boldsymbol{\delta}[j]$, so that the attraction law commands $-k(\mathbf{e} - \boldsymbol{\delta})$, where $\mathbf{e}$ is the centroid error of Section~II-B. With the formation held, and apart from friction and the cap, the centroid velocity follows the frictionless update $\tilde{\mathbf{v}}$ of (\ref{eq:momentum}). The $x$ and $y$ axes then decouple, and for each axis $a$ the state $\mathbf{x} = [e_a,\ \dot{p}_{c,a}]^T$ evolves as
\begin{equation}\label{eq:cl}
\mathbf{x}[j{+}1] = \boldsymbol{\Phi}\,\mathbf{x}[j] + \boldsymbol{\Gamma}\,\delta_a[j],
\end{equation}
where
\[
\boldsymbol{\Phi} = \begin{bmatrix} 1{-}(1{-}\alpha)k\Delta t & \alpha\Delta t \\ -(1{-}\alpha)k & \alpha \end{bmatrix}, \quad \boldsymbol{\Gamma} = (1{-}\alpha)k\begin{bmatrix}\Delta t\\ 1\end{bmatrix}.
\]
The Jury test on the characteristic polynomial $z^2 - (1 + \alpha - (1-\alpha)k\Delta t)z + \alpha$ of $\boldsymbol{\Phi}$ gives exponential stability if and only if
\begin{equation}\label{eq:kmax}
0 < k\Delta t < 2(1+\alpha)/(1-\alpha).
\end{equation}
At $\Delta t = 0.1$~s, $\alpha = 0.70$, and $k = 1$~s$^{-1}$, $k\Delta t = 0.1$ against a limit of \rev{11.3}, and both poles have magnitude $\sqrt{\alpha} = \rev{0.84}$. % [REVIEW] 11.3 = 2(1+alpha)/(1-alpha), 0.84 = sqrt(alpha); gamma from stability_note (outside the authorized source list)
A stable linear loop is input-to-state stable, so a bounded estimate error $\boldsymbol{\delta}$, and friction acting as a command disturbance of size at most $v_c$, leave the centroid error bounded. The speed cap shortens the velocity without turning it, and it is inactive within $v_{\max}/k$ of the estimate, where this linear analysis applies.

\subsubsection{Friction and Estimate Error}
Coulomb friction stops the cluster wherever its command is at most $v_c$. The attraction law commands $k\|\mathbf{e} - \boldsymbol{\delta}\| = k\|\mathbf{p}_c - \hat{\mathbf{p}}^*\|$, so the cluster comes to rest inside the friction ring, a disk of radius $v_c/k$ about the estimate, 0.025~m at $k = 1$~s$^{-1}$. The estimate error shrinks as the cluster nears the critical point (\ref{eq:bound}) and settles at the value it has there. Under ideal continuous kinematics, $\dot{\mathbf{e}} = -k(\mathbf{e} - \boldsymbol{\delta})$, and $V = \frac{1}{2}\|\mathbf{e}\|^2$ satisfies $\dot{V} \le -k\|\mathbf{e}\|(\|\mathbf{e}\| - \|\boldsymbol{\delta}\|)$, so the error decreases while the cluster is farther from the critical point than the estimate is. Provided that moving the cluster a given distance changes the estimate error by less than that distance, $\|\partial\boldsymbol{\delta}/\partial\mathbf{p}_c\| < 1$, the changing estimate causes neither instability nor continued hunting, and the cluster comes to rest. It rests within $v_c/k$ of the estimate, which lies $\|\boldsymbol{\delta}_{\text{final}}\|$ from the critical point, so $\|\mathbf{e}_{\text{final}}\| \le v_c/k + \|\boldsymbol{\delta}_{\text{final}}\|$.

\subsubsection{Near-Singular Estimates}
The fit matrix $\mathbf{A}$ is singular only if the robots become collinear, which the equilateral formation prevents. The estimate can therefore become singular only through $\hat{\mathbf{J}}$, as in near-uniform flow far from a critical point or on a curve where $\det\mathbf{J} = 0$ between two critical points. When $|\det\hat{\mathbf{J}}|$ falls below a small threshold such as $10^{-10}$, the fit gives no usable estimate, and the controller needs a fallback, such as holding the last valid estimate.
The command limit caps the centroid speed in the meantime. No conditioning gate on $\sigma_{\min}(\hat{\mathbf{J}})$ was applied to (\ref{eq:9}).

\subsubsection{Orbits}
The orbital law (\ref{eq:11}) is undefined at $r = 0$, which the ideal dynamics never reach but a noisy estimate can. When $r < 10^{-6}$~m, the controller could simply take a step in a random direction, which carries the cluster off the singular point on the next cycle.

Lag and friction set the steady orbit radius even with an exact estimate. Friction takes $v_c$ off the orbit speed, $V = k_t - v_c$, and passes only the fraction $V/k_t$ of the radial command. The lag makes the cluster follow the turning command a time $T$ late, so in a steady orbit of radius $r^*$ it drifts outward at $TV^2/r^*$. The radial term of (\ref{eq:11}) cancels this drift at
\begin{equation}\label{eq:orbit_cf}
r^* = \tfrac{1}{2}\bigl(r_d + \sqrt{r_d^2 + 4C}\bigr), \qquad C = TVk_t/k_r.
\end{equation}
In the sampled loop $T = \alpha'\Delta t/(1-\alpha') + \Delta t/2$, where $\alpha' = \alpha V/(\alpha V + (1-\alpha)k_t)$ is the lag pole after friction. Without sampling or friction, $T = \tau$ and $C = \tau V^2/k_r$. The error $r^* - r_d$ approaches the minimum radius $r_{\min} = \sqrt{C}$ as $r_d \to 0$ and falls as $C/r_d$ at large radii.



\section{Control Performance in Simulation}
% 2026-11-04: nonlinear fields merged into IV-A to IV-C, former IV-D cut, former Tables II and III merged into Table II. \newnum{} marks numbers not yet reviewed.

A Python-based simulation implementing the three-layer control architecture described in Section~III was developed to verify the critical point estimation and control primitives in Section~II. Unless stated, all trials used a cluster of three robots in an equilateral triangular formation with a 0.33~m side length, matching the hardware configuration (Section~V), with control gains of $k = 1$~s$^{-1}$ in (\ref{eq:10}), and $k_t = 0.2$~m/s, $k_r = 0.5$~s$^{-1}$ in (\ref{eq:11}). The robots follow the model of Section~III-A, and the commanded centroid speed was limited to 0.5~m/s.

Six noise-free analytical fields, one for each of the canonical critical point types listed in Table~\ref{tab:field_summary}, were used for trials, each with its critical point at the origin: sink $(-x,-y)$, source $(x,y)$, saddle $(-y,-x)$, vortex $(-y,x)$, sinking vortex $0.4(-y,x)-0.15(x,y)$, and spewing vortex $0.4(y,-x)+0.15(x,y)$. After performance in noise-free environments was verified, simulations were performed on two fields reconstructed from robot-sensed measurements collected on the hardware testbed (vortex and saddle). These reconstructed fields are a fixed spatial distortion of the analytical fields, fit to the robots' readings on the printed maps, and they add no noise that varies in time.

Three nonlinear fields test how the estimate extrapolates. One cell of the double gyre, $\mathbf{v} = 0.1\pi(\cos\pi x \sin\pi y,\allowbreak\ -\sin\pi x \cos\pi y)$, has a center at the origin, saddles at $(\pm 0.5, \pm 0.5)$~m, and $\det\mathbf{J} = 0$ on the diamond $|x| + |y| = 0.5$~m. A curved saddle, $(x, -y)$ at the origin, bends the inflow axis to $v = -y + \chi y^3$, in an irrotational form, $u - iv = z + \chi z^3$ with $z = x + iy$, and a rotational form with stream function $S(x)S(y)$, where $S(s) = \sin(\omega s)/\omega$ and $\omega^2 = 6\chi$. At the default $\chi = 1$, $\det\mathbf{J}$ vanishes only at $(0, \pm 0.58)$~m in the irrotational form and on the diamond $|x| + |y| = 0.64$~m in the rotational form. Simulations were therefore run on eleven fields in total.

Simulation and hardware trials are scored with the same metrics.

\emph{Attraction.} A trial ends at time $T$ and succeeds when the critical point lies inside the triangle of the three robots and the centroid never left the map. Its final error $e_f = \|\mathbf{p}_c(T) - \mathbf{p}^*\|$ is reported by its mean $\bar{e}_f$, standard deviation $s_e$, median, 95th percentile (p95), and the share of trials with $e_f$ within 0.03, 0.04, and 0.05~m. The bias is the distance of the mean final position from $\mathbf{p}^*$, and the precision $\sigma$ is the RMS distance of the final positions from that mean, so that bias$^2 + \sigma^2$ is the mean of $e_f^2$. The convergence time $T_{\mathrm{conv}}$ is the median time after which the centroid stays within 0.05~m. Localization error is the mean of $\|\hat{\mathbf{p}}^* - \mathbf{p}^*\|$ over the last 20\% of a trial, and control effort is $E_u = \sum_j \|\mathbf{p}_c[j+1] - \mathbf{p}_c[j]\|^2/\Delta t$. In a field with several critical points, success is also scored at the critical point nearest the final centroid, and the other metrics use only the trials that succeeded at the target $\mathbf{p}^*$.

\emph{Orbit.} An orbit run is scored on its radial error $\|\mathbf{p}_c - \mathbf{p}^*\| - r_d$ after the first 10~s. Its mean is the radial bias, its RMS value the orbit RMSE, and its standard deviation the within-run spread. The standard deviation of run means across repeats is the between-trial spread. A radius is practically achievable when the magnitude of its radial bias is at most $0.25\,r_d$. A run is lost if the centroid leaves the map or turns less than a quarter revolution about $\mathbf{p}^*$ after the first 10~s, and the statistics use the runs that were not lost.

\emph{Classification.} Classification accuracy is the fraction of cycles on which the type of $\hat{\mathbf{J}}$ matches the type of the critical point the trial reached.

\subsection{Convergence to Critical Points}

To test convergence to critical points, the following simulation experiments were performed. For each of the eleven fields, the cluster was initialized at a random position drawn uniformly from a 1~m $\times$ 1~m box centered on the target critical point and the simulation was run for 100 time steps at 10~Hz. This was performed 1000 times per field, from the same starts and headings in every field, and the final centroid position of the cluster in each trial was used to determine the metrics. In the gyre and the rotational saddle, the starts with all three robots inside the $\det\mathbf{J} = 0$ diamond around the target are also scored on their own.

 The cluster converged to the critical point in \rev{100}\% of 1000 trials for every analytical field at every threshold (Table~\ref{tab:simulation_results}). There the estimate is exact, and every trial parks on the friction ring of Section~III-D (Fig.~\ref{fig:sim_convergence}). The reconstructed fields' fixed distortion moves the rest point off that ring and lowers success at 0.03~m, while at least \rev{95}\% of trials still end within 0.05~m. % [REVIEW] revision_metrics/tables/table_II_revised.tex; triangle rule: revision_metrics/results/triangle_*.csv

With lag and friction off, white noise up to 0.03~m/s changes the reconstructed fields' final error by at most \rev{1.7}~mm, so nearly all of their error comes from the fixed distortion. On the analytical fields the covariance (\ref{eq:cov}), propagated through the sampled loop, predicts the closed-loop mean error within \rev{4}\% wherever the field changes across the formation by at least three times the noise. % [REVIEW] hyperparameter study v3 E6: median ratio 1.004, range 0.98 to 1.04, over 66 noise settings

\begin{figure}[!t]
    \centering
    \includegraphics[width=0.95\columnwidth]{fig_sim_convergence_band.pdf}
    \caption{Centroid distance to the critical point in simulated plane-fit attraction, mean and $\pm 1$ standard deviation of ten trials per field that start 0.4~m from the critical point at evenly spaced angles. On the analytical vortex every trial parks on the friction ring $v_c/k$, and the reconstructed fields' distortion moves the rest point off it.} % [REVIEW] experiments/revision_metrics/fig_sim_convergence.py (Coulomb model, seed 0): final mean 0.0250, 0.0308 +- 0.0098, 0.0216 +- 0.0129 m
    \label{fig:sim_convergence}
\end{figure}

\begin{table}[!t]
\scriptsize
\centering
\caption{Simulation results on eleven fields. Brackets give starts with all three robots inside the $\det\mathbf{J} = 0$ diamond, and the accuracy rows use the trials that succeeded. A dash marks a radius at which most orbits were lost, and $^{a}$ one at which \newnum{48}\% were.}
\label{tab:simulation_results}
\setlength{\tabcolsep}{1.3pt}
\rev{\begin{tabular}{l|c|cc|ccc}
\hline
 & \textbf{Anal.} & \multicolumn{2}{c|}{\textbf{Recon.}} & \multicolumn{3}{c}{\textbf{Nonlinear}} \\
 & \textbf{(6)} & \textbf{Vort.} & \textbf{Sadd.} & \textbf{Gyre} & \textbf{Irrot.} & \textbf{Rot.} \\
\hline
\multicolumn{7}{l}{\textit{(a) Attraction}} \\
\hline
Succ.\ triangle (\%) & 100 & 100 & 100 & \newnum{49 [100]} & \newnum{99.9} & \newnum{72 [100]} \\
Succ.\ any c.p.\ (\%) & 100 & 100 & 100 & \newnum{99.9} & \newnum{99.9} & \newnum{100} \\
Bias (m) & 0.0013 & 0.0205 & 0.0127 & \newnum{0.0018} & \newnum{0.0010} & \newnum{0.0012} \\
$\sigma$ (m) & 0.0249 & 0.0247 & 0.0252 & \newnum{0.0262} & \newnum{0.0242} & \newnum{0.0259} \\
$\bar{e}_f$ (m) & 0.0249 & 0.0293 & 0.0262 & \newnum{0.0259} & \newnum{0.0238} & \newnum{0.0258} \\
$s_e$ (m) & 0.0008 & 0.0132 & 0.0104 & \newnum{0.0044} & \newnum{0.0044} & \newnum{0.0027} \\
Median (m) & 0.0250 & 0.0321 & 0.0271 & \newnum{0.0249} & \newnum{0.0231} & \newnum{0.0253} \\
p95 (m) & 0.0250 & 0.0500 & 0.0437 & \newnum{0.0334} & \newnum{0.0309} & \newnum{0.0299} \\
Loc.\ error (m) & $<$0.0001 & 0.0222 & 0.0137 & \newnum{0.0057} & \newnum{0.0066} & \newnum{0.0034} \\
$T_{\mathrm{conv}}$ (s) & 2.1 & 2.3 & 2.3 & \newnum{1.8} & \newnum{2.1} & \newnum{1.9} \\
$E_u$ (m$^2$/s) & 0.075 & 0.071 & 0.072 & \newnum{0.047} & \newnum{0.082} & \newnum{0.063} \\
Succ.\ 3/4/5~cm (\%) & 100 & 45/79/95 & 62/89/100 & \newnum{76/100/100} & \newnum{88/100/100} & \newnum{97/100/100} \\
\hline
\multicolumn{7}{l}{\textit{(b) Orbit radial bias (m), by $r_d$ (m)}} \\
\hline
\quad 0.01 & +0.129 & +0.133 & +0.131 & \newnum{+0.127} & \newnum{+0.129} & \newnum{+0.128} \\
\quad 0.10 & +0.093 & +0.094 & +0.089 & \newnum{+0.080} & \newnum{+0.092} & \newnum{+0.085} \\
\quad 0.20 & +0.067 & +0.089 & +0.067 & \newnum{+0.023} & \newnum{+0.067} & \newnum{+0.038} \\
\quad 0.30 & +0.051 & +0.087 & +0.063 & \newnum{$-$0.044} & \newnum{+0.064$^{a}$} & \newnum{$-$0.017} \\
\quad 0.40 & +0.040 & +0.063 & +0.040 & \newnum{--} & \newnum{--} & \newnum{$-$0.081} \\
\quad 0.50 & +0.033 & +0.026 & +0.001 & \newnum{--} & \newnum{--} & \newnum{--} \\
\hline
\multicolumn{7}{l}{\textit{(c) Classification (\%)}} \\
\hline
No noise & 100 & 91 & 100 & \newnum{100} & \newnum{100} & \newnum{100} \\
Noise 0.02~m/s & 90--100 & 65 & 100 & \newnum{86} & \newnum{100} & \newnum{100} \\
\hline
\end{tabular}} % [REVIEW] experiments/revision_metrics/table_merged.py -> tables/table_sim_merged.tex. Eight-field cells: summary.json and triangle_summary.json (former Tables II and III). Nonlinear cells: nonlinear_summary.json (nonlinear.py, Coulomb model, seed 0; noise seed 1)
\end{table}

In the nonlinear fields, curvature showed in the localization error and in which critical point the cluster reached. Trials that reached the target settled as tightly as in the analytical fields, but the estimate stayed \newnum{5.7}~mm from the gyre's center. The field's cubic terms do not cancel over an equilateral triangle, so this error grew roughly as the cube of the formation side, from \newnum{0.16}~mm at 0.10~m to \newnum{25}~mm at 0.50~m. The attraction law equals $-k\hat{\mathbf{J}}^{-1}\hat{\mathbf{v}}(\mathbf{p}_c)$, a Newton step toward the zero of the local fit, and in continuous time such steps cannot cross a curve where $\det\hat{\mathbf{J}} = 0$. The sign of $\det\hat{\mathbf{J}}$ on the first cycle therefore decided the type of critical point reached in all \newnum{2000} gyre and rotational-saddle trials (Fig.~\ref{fig:nonlinear}(a)). Every start with all three robots inside the diamond reached the target. % [REVIEW] nonlinear.py, nonlinear_summary.json: success blocks (loc. error 5.7 mm in the gyre); size block (gyre mean loc. error 0.16, 1.15, 5.72, 24.98 mm at sides 0.10, 0.20, 0.33, 0.50 m); det_hat_sign_mismatch_n = 0 over the 2000 gyre and rotational trials; formation_in_region 154/154 (gyre) and 334/334 (rotational)

\begin{figure}[!t]
    \centering
    \includegraphics[width=0.95\columnwidth]{fig_nonlinear_triangles.pdf}
    \caption{The gyre cell in simulation, with the $\det\mathbf{J} = 0$ diamond dashed and the formation to scale. (a) Attraction starts, blue where the cluster reached the center and orange where it reached a saddle. The orange trial starts with a robot outside the diamond. (b) Orbits at $r_d = 0.3$~m, kept, and 0.4~m, lost.} % [REVIEW] experiments/revision_metrics/fig_nonlinear.py (Coulomb model; trials from nonlinear.py)
    \label{fig:nonlinear}
\end{figure}

\subsection{Orbital Control}

Orbital control was evaluated on the same eleven fields at different commanded radii. For each radius, the cluster was positioned at the desired radius to initialize the run, then the orbital controller (\ref{eq:11}) was used for 600 time steps. The mean radial error and standard deviation were computed after discarding the first 100 time steps (10~s) as a settling transient. Each run was repeated \rev{100} times per field and radius. % [REVIEW] revision_metrics/tables/table_III_revised.tex caption (100 trials per field and radius)


    Table~\ref{tab:simulation_results}(b) presents the orbital tracking results. The six noise-free fields produced identical performance at each radius, with standard deviations below \rev{0.001}~m. This is expected: on an exactly linear noise-free field the three-robot fit recovers the critical point exactly at every step, so the controller input, and hence the trajectory, is independent of the field type. Mean radial errors had a positive bias, and decreased with increasing radius. The reconstructed fields repeat the trend, with error standard deviations that increase with the commanded radius. By the $0.25\,r_d$ criterion the practical range is $r_d \ge \rev{0.3}$~m on the analytical fields and the reconstructed saddle, and $r_d \ge \rev{0.4}$~m on the reconstructed vortex. % [REVIEW] revision_metrics/tables/table_III_revised.tex

In the nonlinear fields, curvature shifted the orbit radius, and orbits were lost once the formation reached a set where $\det\mathbf{J} = 0$. The gyre's speed grows more slowly than in proportion to the distance from its center, so the fit places the center beyond its true position and the cluster orbits inside $r_d$, by \newnum{0.044}~m at $r_d = 0.3$~m. Across $\chi$ from $-1$ to 1, this effect moved the rotational saddle's mean radius by up to \newnum{0.17}~m, while the irrotational form's opposite shifts along its two axes cancelled to within \newnum{0.02}~m. Most gyre orbits from $r_d = 0.4$~m were lost after crossing the diamond (Fig.~\ref{fig:nonlinear}(b)), and about half the irrotational-saddle orbits at 0.3~m were lost past the points $(0, \pm 0.58)$~m. % [REVIEW] nonlinear.py, nonlinear_summary.json orbit and sweep blocks: gyre r_d 0.3 bias -0.044; rotational max shift +0.174 (chi -0.5, r_d 0.4, against +0.040); irrotational max 0.019 (chi -1, r_d 0.3, 59% kept); gyre lost 62% and 93% at 0.4 and 0.5 m (all crossed the diamond); irrotational chi = 1 lost 48% at 0.3 m

\subsection{Classification}

A center ($\operatorname{tr}\mathbf{J} = 0$) and the canonical sink and source ($\mathbf{J} = \mp\mathbf{I}$, with repeated eigenvalues) each lie on a boundary between types in Table~\ref{tab:field_summary}, and any error in $\hat{\mathbf{J}}$ can move them across it. The classifier therefore counts an eigenvalue pair within about $6^\circ$ of the imaginary or real axis as a center or a node. Table~\ref{tab:simulation_results}(c) gives its accuracy over the last 2~s of each attraction trial. Without noise every field except the reconstructed vortex is labeled correctly on every cycle. With white measurement noise of 0.02~m/s on every reading, the center is labeled correctly on \rev{90}\% of cycles and the star nodes on \rev{95}\%, against \rev{100}\% for the saddle. In the nonlinear fields, the type of whichever critical point the cluster reached was labeled correctly on every cycle without noise, and the gyre's center on \newnum{86}\% of cycles with noise. % [REVIEW] revision_metrics/tables/table_classification_revised.tex, error note; nonlinear: nonlinear_summary.json success blocks cls_target, cls_target_noise; trials ending at another critical point any_cls_acc = 1.0

\section{Hardware Validation}


\subsection{Hardware Testbed and Field Representation}

The testbed used for validation is SCU's Robotic Systems Laboratory Decabot testbed \cite{52}, which has been used to validate multirobot adaptive navigation algorithms for both cluster control and swarm particle optimization techniques \cite{8,53}. Fig.~\ref{fig:fullbed} shows the complete testbed.

\begin{figure}[!t]
    \centering
    \includegraphics[width=0.59\columnwidth]{testbed_with_four.png}
    \caption{Experimental testbed: 1.6 m × 1.6 m printed vector field map, Decabot rovers, and OptiTrack system. Four Decabots are shown; the experiments here used three of this fleet.}
    \label{fig:fullbed}
    \end{figure}

The experimental platform used three omnidirectional Decabot rovers operating within a 1.6~m $\times$ 1.6~m workspace. Each Decabot features three omnidirectional wheels providing holonomic motion capability and a downward-facing TCS34725 RGB color sensor enclosed in a light-shrouded housing to minimize external lighting effects. The robots' RGB sensors are calibrated against a printed palette spanning hue $0$ to $2\pi$ and saturation $0$ to $1$ at 24 intervals in each dimension. Each robot records raw RGB and clear-channel counts over every palette cell. The readings are normalized and stored as a four-dimensional feature vector $(R,G,B,K)$. Two feedforward networks, each with two hyperbolic-tangent hidden layers, are trained to recover the target hue and saturation from this vector. The saturation network has a single output, while the hue network outputs the sine and cosine of hue to respect its circularity. Training pools data across robots so that a single model transfers to an uncalibrated unit without per-robot tuning. Calibration achieved hue RMSE of 0.17 radians (9.7$^\circ$, R$^2$=0.991) and saturation RMSE of 4.4\% (R$^2$=0.951) across the operational range. The encoding and calibration palette follows~\cite{22}. All robots were calibrated after receiving preventive maintenance and servicing for this study.

The adaptive navigation and cluster space controller layers of Section~III run in MATLAB/Simulink on a central computer that communicates with the robots over Wi-Fi at 10~Hz: at each cycle it receives OptiTrack positions and RGB sensor readings, computes the control laws, and returns velocity commands, which each robot's onboard computer (Intel Edison) converts into motor commands. The hardware controller also limited the commanded centroid velocity to 0.3~m/s per axis, a safeguard that reserves actuator capacity for formation keeping.

Vector fields are physically realized as printed floor maps using HSV color encoding, where hue (0--2$\pi$) represents flow direction and saturation (0--1) encodes magnitude. Printed magnitudes are scaled to the 0.3--1.0 saturation range to avoid undersaturated areas and ensure reliable sensor readings, then remapped to 0--1 in software, with saturation 1.0 equal to a field speed of 0.3~m/s. Together with the sensor shroud and self-illumination, this encoding is minimally affected by overhead lighting variations.

This method of representing vector fields in printed floor maps has imperfections in both printing and sensing. Fig.~\ref{fig:measurement_error} shows the analytical vortex and saddle fields, reconstructions of each field from sensed robot measurements, and the reconstruction error (right column). These reconstructed fields were used to fit multilayer perceptron (MLP) interpolation models, referenced in Section~IV as a fixed distortion model for simulation.


\begin{figure}[!t]
    \centering
    \includegraphics[width=0.95\columnwidth]{fig_reconstruction_3x2.png}
    \caption{Analytical fields, sensor-based reconstructions, and reconstruction error for vortex (top) and saddle (bottom) fields.} % submitted measurement_error_comparison.png redrawn 3x2 from the same robot readings
    \label{fig:measurement_error}
\end{figure}


\subsection{Validation Test Plan}



Hardware validation focused on two representative field types that encompass the key challenges of vector field navigation: rotational dynamics (vortex) and hyperbolic instability (saddle). The vortex field tests pure rotational dynamics and provides continuity with previous studies in our research group~\cite{22,48}. The saddle is the only canonical field type that contains separatrices, which act as Lagrangian coherent structures in steady flows.

For each field, eight starting locations were selected around the workspace perimeter and the attraction controller (\ref{eq:10}) was run 10 times from each location for \open{20}~s using the three-robot equilateral triangle formation with a \open{0.33 or 0.30}~m side, yielding 80 trials per field. Run order was randomized. Every trial is scored with the metrics of Section~IV, and its final centroid speed is reported so that a cluster still moving is not counted as settled. The classifier of Section~IV-C labels $\hat{\mathbf{J}}$ on every cycle without feeding the controller, and its accuracy is scored against the type of the printed field. % [OPEN] run length, formation side (CAMPAIGN_PROMPT 4.2, 6.2)

Additionally, each field was used to test the critical point orbiting controller (\ref{eq:11}). A total of \open{six} commanded radii were chosen to match those in simulation, and each was run \open{three} times. Each orbiting controller experiment was run for \open{1800} cycles at 10~Hz, allowing the cluster to complete several orbits around the critical point. % [OPEN] number of radii, repeats, orbit run length (CAMPAIGN_PROMPT 6.2)

 \subsection{Attraction Controller Results}

Following the test plan above, \syn{80} vortex trials and \syn{78} saddle trials were analyzed. The final centroid position of each trial was used to determine systematic bias and precision, matching the simulation metrics in Section~IV. % [SYNTHETIC] run counts of the old 3R logs, table_IV_revised.tex; replaced by the campaign

The cluster converged to the critical point in \syn{100}\% of the saddle trials and \syn{86}\% of the vortex trials at the 0.05~m threshold, and in \syn{77}\% and \syn{34}\% at 0.03~m. No trial left the map, and the critical point ended inside the robots' triangle in \syn{100}\% of trials. Table~\ref{tab:estimation_results} summarizes the results. In the vortex field, systematic bias was \syn{0.018}~m with precision of \syn{0.036}~m. In the saddle field, systematic bias was \syn{0.007}~m with precision of \syn{0.024}~m. The bias is smaller than the mean final error because final positions on opposite sides of $\mathbf{p}^*$ cancel in the mean, and the estimate ended closer to the critical point than the cluster. The type of $\hat{\mathbf{J}}$ matched the printed field on \open{campaign}\% of cycles in the vortex and \open{campaign}\% in the saddle. % [SYNTHETIC] old 3R logs scored by revision_metrics (table_IV_revised.tex, V_VIA numbers.tex), replaced by the campaign; [OPEN] hardware classification

Fig.~\ref{fig:vortextomiddle} shows convergence trajectories from all eight starting positions and mean centroid distance to the critical point over time for the same trials.

% [PLACEHOLDER] submitted Figs. 8 and 9, reused until the campaign. (a, b): allrunsadjusted.png split into
% its two panels; the vortex panel shows four-robot runs. (c): fig9.png with its title cropped; its vortex
% curve is the four-robot set (n=79) and its distances are in map units (OptiTrack / 1.6764).
% Replace all three panels with campaign data in metres.
\begin{figure}[!t]
    \centering
    \begin{minipage}[b]{0.45\columnwidth}\centering\includegraphics[width=\textwidth]{placeholder_hw_trajectories_a.png}\end{minipage}\hspace{0.04\columnwidth}%
    \begin{minipage}[b]{0.45\columnwidth}\centering\includegraphics[width=\textwidth]{placeholder_hw_trajectories_b.png}\end{minipage}\\[2pt]
    \includegraphics[width=0.87\columnwidth]{placeholder_hw_distance.png}\\[-1pt]
    {\footnotesize (c)}
    \caption{\syn{[Placeholder]} Convergence trajectories from eight starting positions to the center of the vortex (a) and saddle (b), with the solid line showing the mean path and the shaded region one standard deviation across the 10 repetitions. (c) Distance to center over time, with mean and standard deviation bounds.}
    \label{fig:vortextomiddle}\label{fig:Tomiddlewithtime}
\end{figure}

\begin{table}[!t]
    \footnotesize
    \centering
    \caption{Hardware experimental results \syn{(placeholder data)}. (a) Attraction, vortex: \syn{80} runs, saddle: \syn{78} runs. Bias measures mean offset from origin; precision ($\sigma$) measures repeatability. (b) Orbital control, radial errors in metres after the first 10~s. $\bar e_r$ is the radial bias, $s_b$ the between-trial spread, and $s_w$ the within-run standard deviation. Negative error indicates the cluster orbited inside the commanded radius.} % [SYNTHETIC] table_IV_revised.tex run counts (old 3R logs); V_VIA results/tab_orbit_body.tex
    \label{tab:estimation_results}\label{tab:orbital_radius_comparison}
    \setlength{\tabcolsep}{4pt}
    \syn{\begin{tabular}{l|ccc|ccc}
    \hline
     & \multicolumn{3}{c|}{\textbf{Vortex (printed map)}} & \multicolumn{3}{c}{\textbf{Saddle (printed map)}} \\
    \hline
    \multicolumn{7}{l}{\textit{(a) Attraction}} \\
    \hline
    Bias (m) & \multicolumn{3}{c|}{0.018} & \multicolumn{3}{c}{0.007} \\
    $\sigma$ (m) & \multicolumn{3}{c|}{0.036} & \multicolumn{3}{c}{0.024} \\
    $\bar{e}_f \pm s_e$ (m) & \multicolumn{3}{c|}{0.037 $\pm$ 0.017} & \multicolumn{3}{c}{0.023 $\pm$ 0.009} \\
    Median [p95] (m) & \multicolumn{3}{c|}{0.036 [0.070]} & \multicolumn{3}{c}{0.024 [0.037]} \\
    Loc.\ error (m) & \multicolumn{3}{c|}{0.028} & \multicolumn{3}{c}{0.018} \\
    $T_{\mathrm{conv}}$ (s) & \multicolumn{3}{c|}{3.6} & \multicolumn{3}{c}{3.3} \\
    $E_u$ (m$^2$/s) & \multicolumn{3}{c|}{0.187} & \multicolumn{3}{c}{0.205} \\
    Succ.\ 3/4/5 cm (\%) & \multicolumn{3}{c|}{34/60/86} & \multicolumn{3}{c}{77/99/100} \\
    \hline
    \multicolumn{7}{l}{\textit{(b) Orbit}} \\
    \hline
    $r_d$ (m) & $\bar e_r$ & $s_b$ & $s_w$ & $\bar e_r$ & $s_b$ & $s_w$ \\
    \hline
    0.01 & +0.166 & 0.031 & 0.027 & +0.111 & 0.014 & 0.018 \\
    0.10 & +0.151 & 0.028 & 0.061 & +0.074 & 0.015 & 0.028 \\
    0.20 & +0.128 & 0.027 & 0.064 & +0.071 & 0.018 & 0.047 \\
    0.30 & +0.112 & 0.027 & 0.097 & +0.091 & 0.033 & 0.059 \\
    0.40 & +0.081 & 0.036 & 0.119 & +0.102 & 0.040 & 0.079 \\
    0.50 & +0.018 & 0.020 & 0.079 & +0.088 & 0.041 & 0.095 \\
    \hline
    \end{tabular}} % [SYNTHETIC] revision_metrics/tables/table_IV_revised.tex (old 3R logs) and V_VIA results/tab_orbit_body.tex (three-robot columns, rows 0.05 and 0.15 m omitted), replaced by the campaign
    \end{table}



\subsection{Orbital Control Results}

Orbital control was evaluated on both fields at commanded radii from 0.01~m to 0.50~m. For each radius, the orbital controller (\ref{eq:11}) was run for \open{1800} cycles at 10~Hz, and the mean radial error and its standard deviation were recorded after the first 10~s, as in simulation (Table~\ref{tab:orbital_radius_comparison}). % [OPEN] orbit run length (CAMPAIGN_PROMPT 6.2)

 In the vortex field, mean radial errors decreased from \syn{0.166}~m at a commanded radius of 0.01~m to \syn{0.018}~m at 0.50~m, with within-run standard deviations ranging from \syn{0.027} to \syn{0.119}~m. In the saddle field, mean radial errors ranged from \syn{0.111}~m to \syn{0.088}~m over the same commanded radii, with within-run standard deviations increasing from \syn{0.018} to \syn{0.095}~m. From $r_d = \syn{0.10}$~m upward the fluctuation within one orbit exceeds the variation between repeats. By the $0.25\,r_d$ criterion the practical range is $r_d \ge \syn{0.40}$~m in the vortex and $r_d \ge \syn{0.50}$~m in the saddle. Below these radii the controller holds a bounded orbit about the critical point but does not track the commanded radius. \syn{One} vortex and \syn{three} saddle orbits failed by leaving the map, all at $r_d \ge \syn{0.40}$~m. % [SYNTHETIC] V_VIA numbers.tex and tab_orbit_body.tex (three-robot columns, six radii)




\section{Discussion}

\subsection{Error Analysis}

Estimation and control errors arise from four sources: robot dynamics, field distortion and measurement noise, cluster geometry, and field conditioning. Sections~II-B and~III-D bound the effect of each, and this section traces the errors of Sections~IV and~V to them.

Perfect convergence does not occur in simulation. On the analytical fields the estimate $\hat{\mathbf{p}}^*$ matched $\mathbf{p}^*$ to machine precision, yet every trial stopped on the friction ring of radius $v_c/k$ (Section~III-D), where the attraction gain $k$ times the remaining distance no longer overcomes the friction level $v_c$. Some errors may also have arisen from modeling discrepancies such as modeling all robots identically with the same friction level, speed limit, and time constant. % [REVIEW] revision_metrics/results/summary.json: analytical est_ss_mean 2e-18 m (Table II Loc. error < 0.0001)

Positive radial bias at small orbital radii is also a result of robot dynamics. At small radii, the commanded angular rate $k_t / r_d$ is high relative to the radial correction bandwidth $k_r$, and the outward drift from the actuator lag sets the bias (\ref{eq:orbit_cf}). The closed form matches simulation to a median of \rev{0.03}~mm across sweeps of gains, lag, sample period, and friction. At the settings of Section~IV, $r_{\min} = \rev{0.133}$~m, and the error falls below 0.05~m only for $r_d \ge \rev{0.31}$~m. Lowering the orbital speed $k_t$ to 0.1~m/s cuts $r_{\min}$ to \rev{0.058}~m, and removing the lag cuts it to \rev{0.059}~m. Raising the radial gain $k_r$ to 2~s$^{-1}$ cuts it to \rev{0.067}~m, at the cost of increased sensitivity to noise in the critical point estimate. % [REVIEW] hyperparameter study v3 report_numbers.txt (closed form, table 3); C = T V k_t / k_r is theory.orbit_C

Comparing the noise-free and reconstructed fields in simulation isolates the effect of field distortion: the fixed distortion of the printed maps appears to be the driving factor of the simulation-reality gap, with sensing errors propagating directly through the critical point estimation equations. Noise that varies in time is far smaller. With the formation parked at fixed points, a robot's reading changes by about \syn{$2\times10^{-5}$}~m/s between cycles \open{campaign block E}, while the calibration error of Section~V-A is about 0.02~m/s and fixed at each location. In the final error of Section~III-D, friction sets the spread of final positions through $v_c/k$, and the distortion sets their bias through the estimate error at the critical point. % [SYNTHETIC] old hardware logs (prose batch N1); campaign block E (static noise runs) replaces it

The hardware orbit metrics separate estimation error from control error. The controller-only bias measures the radius about $\hat{\mathbf{p}}^*$, which (\ref{eq:11}) regulates, instead of about $\mathbf{p}^*$. At $r_d = 0.01$~m it was \syn{0.177}~m in the vortex against a true bias of \syn{0.166}~m, so the small-radius offset arises in the closed loop. At larger radii the localization error grows with distance from the critical point, and the within-run spread grows with it. % [SYNTHETIC] V_VIA numbers.tex

To avoid errors arising from robots becoming collinear (the condition number of $\mathbf{A}$ becomes infinite), an equilateral triangle was chosen as the cluster shape. This shape minimizes the condition number of $\mathbf{A}$ in (\ref{eq:4}) among formations of the same size centered on the sampled region, making this formation optimal for sampling.

Poor conditioning of the estimated Jacobian amplifies coefficient errors into errors in the critical point estimate (\ref{eq:9}) by the factor $1/\sigma_{\min}(\hat{\mathbf{J}})$ in (\ref{eq:bound}). Unlike the condition number of $\mathbf{A}$, it is a property of the field and is not under the designer's control. Where $\det\hat{\mathbf{J}}$ approaches zero, so does $\sigma_{\min}(\hat{\mathbf{J}})$, and Section~III-D describes how the controller handles it. A center is non-hyperbolic yet has $\det\mathbf{J} > 0$, so (\ref{eq:9}) locates it as well as any other critical point. Its type is harder to label, because its eigenvalues lie on the boundary between the two spiral types (Section~IV-C).

The measurement term of (\ref{eq:bound}) grows with the extrapolation distance $\|\mathbf{e}\|/R$, and the reconstructed field degrades near the edge of the mapped region. Together they account for the orbital bias at the largest commanded radius falling below the controller-only bias on both reconstructed fields, to \rev{+0.026} against \rev{+0.035}~m on the vortex and \rev{+0.001} against \rev{+0.017}~m on the saddle. There the estimated distance to the critical point is inflated, and the radial controller pulls the cluster inward. % [REVIEW] revision_metrics/results/summary.json, reconstructed orbit r_d 0.50 (bias_w1, ctrl_bias_w1)

\subsection{Comparison with the Vector-to-Scalar Primitive}

The modular architecture allows navigation controllers to be interchanged while cluster and robot controller layers remain fixed. As a result, performance differences between controllers reflect algorithmic rather than implementation variations. % [S] both sentences
The vector-to-scalar primitive \cite{48} reduces the flow to its magnitude, the speed $S_i = \|\mathbf{v}(\mathbf{p}_i)\|$ at each robot, and a plane fit to the three speeds gives their gradient $\mathbf{g}$. For attraction the cluster descends $\mathbf{g}$. For orbits the primitive was modified to move tangentially to the direction of steepest descent. It was run in simulation with the robot model, gains, and formation of Section~IV in three forms (Table~\ref{tab:baselines}). The first is the primitive alone. The second, the oracle, takes the true critical point as the target of (\ref{eq:10}) and the center of the radial term of (\ref{eq:11}). The third places the critical point where the fitted speed reaches zero, $\hat{\mathbf{p}}_m = \mathbf{p}_c - S_c\mathbf{g}/\|\mathbf{g}\|^2$ with $S_c$ the fitted speed at the centroid, and feeds it to (\ref{eq:10}) and (\ref{eq:11}).

On the six analytical fields the speed grows equally in every direction, so its gradient points directly away from $\mathbf{p}^*$, and the primitive reached $\mathbf{p}^*$ in every trial (Fig.~\ref{fig:effort}(a)). The fitted gradient keeps the same length at every distance, so the command holds the robots at their speed limit until the cluster passes $\mathbf{p}^*$, overshoots by a median of \newnum{0.031}~m, and rebounds. The Jacobian estimate scales the command with distance, and the cluster slows along a smooth exponential. The primitive settles faster and closer to $\mathbf{p}^*$, because friction stops a proportional command on the ring $v_c/k$, but at \newnum{1.9} times the control effort. % [REVIEW] baselines_v2s_summary.json linear: v2s overshoot_med 0.031, e_mean 0.003, t_conv 1.1, E_u 0.137 against 0.074

In orbit, the primitive alone has no mechanism to detect or correct radial error with respect to a specific point, and so centrifugal drift goes uncorrected. Every orbit spiraled outward off the map (Fig.~\ref{fig:effort}(b)). Given the oracle, the same tangent with the radial term of (\ref{eq:11}) matched the Jacobian method within \newnum{0.002}~m RMSE at every radius. With its own speed estimate the primitive held the radius at a lower RMSE, because the estimate's inward bias offsets the lag drift of Section~III-D. That estimate never lets the cluster rest during attraction, however, and the cluster still moved at \newnum{0.12}~m/s after 10~s. % [REVIEW] baselines_v2s_summary.json linear: v2s lost 100% at every r_d (all left the map); v2s_oracle RMSE 0.092/0.052/0.035 against 0.093/0.051/0.033; v2s_est RMSE 0.068/0.016/0.020, final_speed 0.116, E_u 0.306

The speed gradient points away from the critical point only where the speed grows equally in every direction. The anisotropic node $\mathbf{J} = \operatorname{diag}(-1, -0.2)$ is affine, but its speed grows five times more slowly along one axis. There the primitive reached $\mathbf{p}^*$ in \newnum{7}\% of trials, and most of the rest ended on the slow axis, a median \newnum{0.67}~m from $\mathbf{p}^*$. Its speed estimate erred by \newnum{0.46}~m on average during orbits, and even with the oracle most orbits were lost, because the speed contours are ellipses. The Jacobian estimate is exact on any affine field, including the node. In the gyre, the primitive reached the same critical point as the Jacobian method in \newnum{96}\% of trials. The speed vanishes at the center and at the saddles alike, however, so it cannot tell which it has reached, while the Jacobian method labeled the type on every cycle (Section~IV-C). The primitive lost every gyre orbit, and with the oracle every orbit from $r_d = 0.3$~m, where the gyre's speed contours no longer circle the center. % [REVIEW] baselines_v2s_summary.json node_aniso: v2s success 6.6% (467 failures, median |y| 0.67 m, |x| 0.03 m), v2s_est est_orb 0.455, v2s_oracle lost 90/82/100%; gyre: same point as jacobian 96.4%, jacobian type_acc 100%, v2s lost 100%, v2s_oracle lost 0/100/100%, jacobian lost 0/0/100%

\begin{figure}[!t]
    \centering
    \includegraphics[width=0.95\columnwidth]{fig_v2s_effort.pdf}
    \caption{Simulated attraction and orbit on the analytical vortex with the Jacobian estimate and the vector-to-scalar (V2S) primitive, mean and $\pm 1$ standard deviation of ten runs from evenly spaced starts. (a) Centroid speed from 0.4~m, with control effort $E_u$ in the legend. (b) Distance from the critical point at the commanded $r_d = 0.3$~m (dashed).} % [REVIEW] experiments/revision_metrics/fig_v2s_effort.py, draft A (Coulomb model, seed 0)
    \label{fig:effort}
\end{figure}

\begin{table}[!t]
\centering
\scriptsize
\caption{Comparison with the vector-to-scalar (V2S) primitive in simulation, \newnum{500} attraction trials of 10~s and \newnum{50} orbit runs of 60~s per field and setting. The accuracy rows use the trials that succeeded, gyre success counts the center, and a dash marks a radius at which most orbits were lost.}
\label{tab:baselines}
\setlength{\tabcolsep}{1.8pt}
\newnum{\begin{tabular}{ll|cccc|ccc}
\hline
 & & \multicolumn{4}{c|}{\textbf{Attraction}} & \multicolumn{3}{c}{\textbf{Orbit RMSE (m)}} \\
\textbf{Field} & \textbf{Navigation} & \textbf{Succ.} & \textbf{$\bar{e}_f$} & \textbf{$T_{\mathrm{conv}}$} & \textbf{$E_u$} & \multicolumn{3}{c}{\textbf{at $r_d$ (m)}} \\
 & & \textbf{(\%)} & \textbf{(m)} & \textbf{(s)} & \textbf{(m$^2$/s)} & 0.1 & 0.3 & 0.5 \\
\hline
Linear (6) & Jacobian & 100 & 0.025 & 2.1 & 0.074 & 0.093 & 0.051 & 0.033 \\
 & V2S & 100 & 0.003 & 1.1 & 0.137 & -- & -- & -- \\
 & V2S + oracle & 100 & 0.025 & 2.1 & 0.074 & 0.092 & 0.052 & 0.035 \\
 & V2S + speed est. & 100 & 0.012 & 1.3 & 0.306 & 0.068 & 0.016 & 0.020 \\
\hline
Aniso.\ node & Jacobian & 100 & 0.025 & 2.1 & 0.074 & 0.093 & 0.051 & 0.033 \\
 & V2S & 7 & 0.063 & 6.8 & 0.103 & -- & -- & -- \\
 & V2S + oracle & 100 & 0.025 & 2.1 & 0.074 & -- & -- & -- \\
 & V2S + speed est. & 6 & 0.071 & 10.0 & 0.228 & -- & -- & -- \\
\hline
Gyre & Jacobian & 50 & 0.026 & 1.8 & 0.044 & 0.080 & 0.049 & -- \\
 & V2S & 48 & 0.009 & 0.9 & 0.085 & -- & -- & -- \\
 & V2S + oracle & 100 & 0.025 & 2.1 & 0.074 & 0.091 & -- & -- \\
 & V2S + speed est. & 48 & 0.013 & 1.0 & 0.273 & 0.068 & -- & -- \\
\hline
\end{tabular}} % [REVIEW] experiments/revision_metrics/baselines_v2s.py --table -> tables/table_v2s.tex
\end{table}

\subsection{Limitations}

Several limitations constrain the current implementation and experimental validation. First, the approach is restricted to planar (2D) vector fields, though many real-world applications involve three-dimensional flows. Extension to 3D fields would require at minimum four non-coplanar robots for Jacobian estimation.

Second, the testbed validates performance only on static fields where critical points remain fixed during operation. The static printed field approach enables precise ground truth validation but does not capture dynamic field behavior, and the printed fields do not push the robots. A physical flow advects the cluster and adds its own velocity $\mathbf{J}\mathbf{e}$ to the centroid near $\mathbf{p}^*$. For an ideal centroid with an exact estimate, the error then obeys $\dot{\mathbf{e}} = (\mathbf{J} - k\mathbf{I})\mathbf{e}$, which is unstable when an eigenvalue of $\mathbf{J}$ has real part above $k$. A physical critical point can also drift, and the attraction law trails one moving at constant speed by that speed divided by $k$.

Third, robot team sizes were limited to three. While the theoretical framework supports arbitrary team sizes through the overdetermined system approach, larger formations present practical implementation challenges including communication bandwidth, computational complexity, and formation control stability. With $n > 3$ robots the least-squares fit averages measurement errors across robots and survives the loss of one. However, a wider footprint asks the affine model to hold over more of the field's curvature.

Fourth, the estimator assumes the field is well approximated by an affine model from the cluster out to the critical point (Section~II-B). On fields that are highly nonlinear at the scale of the formation, this local linear fit degrades, biasing the recovered critical point location. Section~IV probes this regime on a gyre cell and a curved saddle, but fields with stronger curvature, or with critical points closer together than one formation side, were not tested.



Finally, the computation is centralized. A central computer receives OptiTrack positions and field readings from every robot, forms the estimate, and returns velocity commands. The estimate needs only the positions and readings of one cycle in a common frame, so a robot receiving them from its neighbors could compute it. Onboard relative localization and inter-robot communication were not tested. A dropout can leave fewer than three readings, and the fit then has no unique solution, while latency mixes readings from different times into one fit. The robustness of the approach to lower update rates and communication dropouts was not tested.

\subsection{Future Work}

The estimation framework uses only instantaneous measurements and carries no state between control cycles. This memoryless property means the estimator does not assume field stationarity, suggesting applicability to environments where the field itself varies with time. Validating the estimator on a genuinely time-varying field remains future work. Beyond critical point detection and orbital control, vector fields contain other topological features of practical interest. Separatrices connected to saddle points define boundaries between distinct flow regimes and act as Lagrangian coherent structures in steady flows. The three-robot fit returns a spatially constant Jacobian, so detecting them calls for a second-order fit, which six robots that do not lie on a common conic can provide.



\section{Conclusion}
\noindent

This paper presents a framework for multirobot navigation in vector fields containing critical points, in which the robots sense the field at spatially distributed points and a central computer forms the estimate. A Jacobian-based approach uses instantaneous measurements from three robots to identify critical point locations without prior field knowledge, supporting attraction and orbital control laws built on these estimates. Simulation across six canonical field types confirmed convergence of these control primitives. On curved fields the estimate degraded as the cube of the formation side, and the sign of $\det\hat{\mathbf{J}}$ on the first cycle decided which type of critical point the cluster reached. In \syn{158} hardware convergence trials, the system achieved a \syn{100}\% (saddle) and \syn{86}\% (vortex) success rate within 0.05~m, with mean position errors of \syn{0.037}~m or less. Orbital control experiments maintained commanded radii within a quarter of the radius from \syn{0.40}~m (vortex) and \syn{0.50}~m (saddle) upward, and below that held bounded orbits that did not track the commanded radius. % [SYNTHETIC] hardware numbers
The printed fields do not advect the robots, so these results hold under lab conditions. A physical flow adds a disturbance to the robot dynamics that the present controller does not compensate for.


\section*{Disclosure of Conflicts of Interest}

C. A. Kitts serves as an Associate Editor of the \textit{IEEE Systems Journal}. He has been excluded from all editorial handling and peer review of this article, including the assignment of an Associate Editor and the selection of reviewers, and has been designated a non-preferred Associate Editor in the submission system.

C. Waight and C. A. Kitts are co-inventors on a provisional patent application filed through Santa Clara University covering the method described in this article. Santa Clara University holds any rights arising from that application.

The authors declare no other financial or non-financial competing interests.


%\balance % off while the last page has room; restore once the enlarged figures fill it
\begin{thebibliography}{99}

\bibitem{1} C. Kitts and M. Egerstedt, "Design, control, and applications of real-world multirobot systems [from the guest editors]," \textit{IEEE Robot. Autom. Mag.}, vol. 15, no. 1, p. 8, Mar. 2008.

\bibitem{2} J. Cortés and M. Egerstedt, "Coordinated control of multi-robot systems: A survey," \textit{SICE J. Control Meas. Syst. Integr.}, vol. 10, no. 6, pp. 495--503, 2017.

\bibitem{3} P. Ögren, E. Fiorelli, and N. E. Leonard, "Cooperative control of mobile sensor networks: Adaptive gradient climbing in a distributed environment," \textit{IEEE Trans. Autom. Control}, vol. 49, no. 8, pp. 1292--1302, 2004.

\bibitem{4} L. Briñón-Arranz, A. Renzaglia, and L. Schenato, "Multirobot symmetric formations for gradient and hessian estimation with application to source seeking," \textit{IEEE Trans. Robot.}, vol. 35, no. 3, pp. 782--789, 2019.

\bibitem{5} C. A. Kitts, R. T. McDonald, and M. A. Neumann, "Adaptive navigation control primitives for multirobot clusters: Extrema finding, contour following, ridge/trench following, and saddle point station keeping," \textit{IEEE Access}, vol. 6, pp. 17625--17642, 2018.

\bibitem{6} R. T. McDonald, C. A. Kitts, and M. A. Neumann, "Experimental implementation and verification of scalar field ridge, trench, and saddle point maneuvers using multirobot adaptive navigation," \textit{IEEE Access}, vol. 7, pp. 62950--62961, 2019.

\bibitem{7} R. K. Lee et al., "Multiple UAV adaptive navigation for three-dimensional scalar fields," \textit{IEEE Access}, vol. 9, pp. 122626--122654, 2021.

\bibitem{8} R. T. McDonald, C. A. Kitts, and M. A. Neumann, "Navigation of scalar fronts with multirobot clusters in simulation and experiment," \textit{IEEE Syst. J.}, vol. 14, no. 3, pp. 3755--3766, 2020.

\bibitem{9} R. K. Lee, C. A. Kitts, M. A. Neumann, and R. T. McDonald, "3-d adaptive navigation: Multirobot formation control for seeking and tracking of a moving source," in \textit{Proc. IEEE Aerosp. Conf.}, pp. 1--12, 2020.

\bibitem{10} T. Adamek, C. A. Kitts, and I. Mas, "Gradient-based cluster space navigation for autonomous surface vessels," \textit{IEEE/ASME Trans. Mechatronics}, vol. 20, no. 2, pp. 506--518, 2014.

\bibitem{11} F. Pagano, N. De Carli, E. Restrepo, A. Marino, and P. Robuffo Giordano, "Distributed multi-robot active-sensing of a diffusive source," \textit{IEEE Robot. Autom. Lett.}, vol. 10, no. 10, pp. 10807--10814, 2025.

\bibitem{13} S. Li, R. Kong, and Y. Guo, "Cooperative distributed source seeking by multiple robots: Algorithms and experiments," \textit{IEEE/ASME Trans. Mechatronics}, vol. 19, no. 6, pp. 1810--1820, 2014.

\bibitem{14} Z. Deng and J. Luo, "Fully distributed algorithms for constrained nonsmooth optimization problems of general linear multiagent systems and their application," \textit{IEEE Trans. Autom. Control}, vol. 69, no. 2, pp. 1377--1384, 2023.

\bibitem{15} W. Li, J. A. Farrell, S. Pang, and R. M. Arrieta, "Moth-inspired chemical plume tracing on an autonomous underwater vehicle," \textit{IEEE Trans. Robot.}, vol. 22, no. 2, pp. 292--307, 2006.

\bibitem{18} S. R. Ramp et al., "Preparing to predict: the second autonomous ocean sampling network (AOSN-II) experiment in the Monterey Bay," \textit{Deep Sea Res. Part II: Top. Stud. Oceanogr.}, vol. 56, no. 3-5, pp. 68--86, 2009.

\bibitem{19} J. Das et al., "Coordinated sampling of dynamic oceanographic features with underwater vehicles and drifters," \textit{Int. J. Robot. Res.}, vol. 31, no. 5, pp. 626--646, 2012.

\bibitem{20} J. Wang et al., "Dynamic plume tracking by cooperative robots," \textit{IEEE/ASME Trans. Mechatronics}, vol. 24, no. 2, pp. 609--620, 2019.

\bibitem{21} K. Garg and D. Panagou, "Finite-time estimation and control for multi-aircraft systems under wind and dynamic obstacles," \textit{J. Guid. Control Dyn.}, vol. 42, no. 7, pp. 1489--1505, 2019.

\bibitem{22} C. Waight and C. Kitts, "A functional indoor testbed for multirobot adaptive navigation in vector field environments," in \textit{Proc. ASME Int. Design Eng. Tech. Conf. Comput. Inf. Eng. Conf.}, vol. 89251, 2025.

\bibitem{23} M. Michini, M. A. Hsieh, E. Forgoston, and I. B. Schwartz, "Robotic tracking of coherent structures in flows," \textit{IEEE Trans. Robot.}, vol. 30, no. 3, pp. 593--603, 2014.

\bibitem{24} D. R. Nelson, D. B. Barber, T. W. McLain, and R. W. Beard, "Vector field path following for miniature air vehicles," \textit{IEEE Trans. Robot.}, vol. 23, no. 3, pp. 519--529, 2007.

\bibitem{25} W. Yao, H. G. de Marina, B. Lin, and M. Cao, "Singularity-free guiding vector field for robot navigation," \textit{IEEE Trans. Robot.}, vol. 37, no. 4, pp. 1206--1221, 2021.

\bibitem{26} D. Panagou, "Motion planning and collision avoidance using navigation vector fields," in \textit{Proc. IEEE Int. Conf. Robot. Autom.}, 2014, pp. 2513--2518.

\bibitem{28} X. Wang et al., "Adaptive vector field guidance without a priori knowledge of course dynamics and wind," \textit{IEEE/ASME Trans. Mechatronics}, vol. 27, no. 6, pp. 4597--4607, 2022.

\bibitem{29} Y. Zhang, M. A. Godin, J. G. Bellingham, and J. P. Ryan, "Using an autonomous underwater vehicle to track a coastal upwelling front," \textit{IEEE J. Ocean. Eng.}, vol. 37, no. 3, pp. 338--347, 2012.

\bibitem{32} D. Kularatne and M. A. Hsieh, "Tracking attracting manifolds in flows," \textit{Auton. Robots}, vol. 41, no. 8, pp. 1575--1588, 2017.

\bibitem{36} X. Tricoche, G. Scheuermann, and H. Hagen, "Continuous topology simplification of planar vector fields," in \textit{Proc. IEEE Vis.}, pp. 159--166, 2001.

\bibitem{37} W. G. Pichel et al., "Marine debris collects within the North Pacific subtropical convergence zone," \textit{Mar. Pollut. Bull.}, vol. 54, no. 8, pp. 1207--1211, 2007.

\bibitem{38} M. Mateus, R. Canelas, L. Pinto, and N. Vaz, "When tragedy strikes: potential contributions from ocean observation to search and rescue operations after drowning accidents," \textit{Front. Mar. Sci.}, vol. 7, p. 55, 2020.

\bibitem{40} P. Keramea, K. Spanoudaki, G. Zodiatis, G. Gikas, and G. Sylaios, "Oil spill modeling: A critical review on current trends, perspectives, and challenges," \textit{J. Mar. Sci. Eng.}, vol. 9, no. 2, p. 181, 2021.

\bibitem{41} J. Wang, Y. Lin, R. Liu, and J. Fu, "Odor source localization of multi-robots with swarm intelligence algorithms: A review," \textit{Front. Neurorobot.}, vol. 16, p. 949888, 2022.

\bibitem{43} C. Dong et al., "Oceanic mesoscale eddies," \textit{Ocean-Land-Atmos. Res.}, vol. 4, p. 0081, 2025.

\bibitem{44} Z. Li, X. Wang, J. Hu, F. P. Andutta, and Z. Liu, "A study on an anticyclonic-cyclonic eddy pair off Fraser Island, Australia," \textit{Front. Mar. Sci.}, vol. 7, p. 594358, 2020.

\bibitem{45} R. Deogharia, H. Gupta, S. Sil, A. Gangopadhyay, and A. Shee, "On the evidence of helico-spiralling recirculation within coherent cores of eddies using Lagrangian approach," \textit{Sci. Rep.}, vol. 14, no. 1, p. 11014, 2024.

\bibitem{30} T. Inanc, S. C. Shadden, and J. E. Marsden, "Optimal trajectory generation in ocean flows," in \textit{Proc. Amer. Control Conf.}, pp. 674--679, 2005.

\bibitem{31} R. N. Smith, M. Schwager, S. L. Smith, B. H. Jones, D. Rus, and G. S. Sukhatme, "Persistent ocean monitoring with underwater gliders: Adapting sampling resolution," \textit{J. Field Robot.}, vol. 28, no. 5, pp. 714--741, 2011.

\bibitem{33} G. Knizhnik, P. Li, X. Yu, and M. A. Hsieh, "Flow-based control of marine robots in gyre-like environments," in \textit{Proc. IEEE Int. Conf. Robot. Autom. (ICRA)}, pp. 3047--3053, 2022.

\bibitem{46} Y. Jiao, H. Hang, J. Merel, and E. Kanso, "Sensing flow gradients is necessary for learning autonomous underwater navigation," \textit{Nat. Commun.}, vol. 16, no. 1, p. 3044, 2025.

\bibitem{47} D. Rajabhandharaks, R. T. McDonald, M. A. Neumann, and C. A. Kitts, "Indoor testbed for vector field multirobot adaptive navigation: Feasibility study," in \textit{Proc. ASME Int. Design Eng. Tech. Conf. Comput. Inf. Eng. Conf.}, vol. 59292, p. V009T12A043, 2019.

\bibitem{48} R. Cooper et al., "Initial Study of Multirobot Adaptive Navigation for Exploring Environmental Vector Fields," in \textit{Proc. IEEE Aerosp. Conf.}, 2019.

\bibitem{49} Z. Zhou, J. Liu, and J. Yu, "A survey of underwater multi-robot systems," \textit{IEEE/CAA J. Autom. Sinica}, vol. 9, no. 1, pp. 1--18, 2021.

\bibitem{50} C. A. Kitts and I. Mas, "Cluster space specification and control of mobile multirobot systems," \textit{IEEE/ASME Trans. Mechatronics}, vol. 14, no. 2, pp. 207--218, 2009.

\bibitem{51} I. Mas and C. Kitts, "Obstacle avoidance policies for cluster space control of nonholonomic multirobot systems," \textit{IEEE/ASME Trans. Mechatronics}, vol. 17, no. 6, pp. 1068--1079, 2012.

\bibitem{52} S. Tomer et al., "A low-cost indoor testbed for multirobot adaptive navigation research," in \textit{Proc. IEEE Aerosp. Conf.}, pp. 1--12, 2018.

\bibitem{53} S. T. Hart and C. A. Kitts, "Unifying control architecture for reactive particle swarms," \textit{IEEE/ASME Trans. Mechatronics}, vol. 28, no. 2, pp. 873--883, 2022.

\end{thebibliography}

\begin{IEEEbiography}[{\includegraphics[width=1in,height=1.25in,clip,keepaspectratio]{fig1.jpg}}]{Christopher Waight}
received his BSc degree in Physics from Santa Clara University. He later earned his Master's in Mechanical Engineering there, and is now a PhD Candidate in Mechanical Engineering. He currently works as a Research Assistant in SCU's Robotic Systems Laboratory. He also has experience working as a hardware engineer and as a Technical Program Manager in R\&D.
\end{IEEEbiography}

\begin{IEEEbiography}[{\includegraphics[width=1in,height=1.25in,clip,keepaspectratio]{ckitts.png}}]{Christopher A. Kitts}
    (Senior Member, IEEE) earned degrees from Princeton University (B.S.E.), University of Colorado (M.P.A.), and Stanford University (M.S. and Ph.D.). He served as a Satellite Constellation Mission Controller with the U.S. Air Force and as a Computer Scientist at NASA Ames Research Center. He is currently Professor of Mechanical Engineering, Director of the Robotic Systems Laboratory, and Associate Dean at Santa Clara University. Dr. Kitts is a Fellow of ASME.
    \end{IEEEbiography}

\end{document}
